% Converted from README.md
\subsection{Jazzer (JVM fuzzer)}\label{jazzer-jvm-fuzzer}

Jazzer omogućava fuzz testiranje Java i drugih JVM programa. U JUnit načinu rada fuzz test se piše kao običan test metod označen anotacijom \texttt{@FuzzTest}. Bez dodatnog podešavanja takav test može da se izvršava kao regresioni test nad sačuvanim ulazima. Kada se postavi promenljiva okruženja \texttt{JAZZER\_FUZZ=1}, Jazzer pokreće generisanje novih ulaza.

Minimalna Maven konfiguracija za primer u ovom direktorijumu koristi \texttt{jazzer-junit} verziju \texttt{0.30.0}.

\codefile{pom.xml}{xml}

Primer ispod proverava jednostavnu šemu kodiranja. Jedinični test proverava jedan poznat ulaz, dok fuzz test proverava isto svojstvo nad velikim brojem ulaza koje bira Jazzer. Anotacije \texttt{@NotNull}, \texttt{@WithUtf8Length} i \texttt{@InRange} ograničavaju prostor ulaza kada je to potrebno za smislen test.

\codefile{src/main/java/org/example/SomeScheme.java}{java}

\codefile{src/test/java/org/example/ParserTests.java}{java}

Testovi se mogu pokrenuti kao obični JUnit testovi:

\begin{minted}{bash}
$ mvn test -Dtest=ParserTests
\end{minted}

Za fuzz režim postavlja se promenljiva okruženja:

\begin{minted}{bash}
$ JAZZER_FUZZ=1 mvn test -Dtest=ParserTests
\end{minted}

Kada Jazzer pronađe ulaz koji ruši proces ili narušava tvrđenje, taj ulaz treba sačuvati kao regresioni test. Na taj način se pronađena greška proverava i posle ispravke.
